% Source PDF: source/普通高考/2014/2014山东理.pdf, page 1, question 11.
% Original PNG reference: img/2014/shandong_science/q11_fig1.png.
% Node -> directed edge list: 开始→输入x→n=0→判定；是→x=x+1→n=n+1→判定；否→输出n→结束。
% The loop returns through the horizontal merge before the unique downward arrow into the decision.
\begin{tikzpicture}[
  x=1cm,y=1cm,
  >=Stealth,
  line width=0.55pt,
  line cap=round,line join=round,
  every node/.style={font=\normalsize,anchor=center,align=center,
    execute at begin node={\setlength{\parindent}{0pt}}},
  flowbox/.style={draw,minimum height=.66cm,inner xsep=4pt,inner ysep=2pt,
    anchor=center,align=center},
  io/.style={draw,shape=trapezium,trapezium left angle=78,trapezium right angle=102,
    trapezium stretches=false,minimum height=.76cm,inner xsep=4pt,inner ysep=2pt,
    anchor=center,align=center},
  flowarrow/.style={->,line width=0.55pt},
]
  \node[draw,rounded corners=7pt,minimum width=1.35cm,minimum height=.70cm]
    (start) at (0,8.55) {开始};
  \node[io,minimum width=1.65cm] (input) at (0,7.25) {输入 $x$};
  \node[flowbox,minimum width=1.45cm] (init) at (0,5.98) {$n=0$};
  \coordinate (merge) at (0,5.43);
  \node[draw,diamond,aspect=2.85,minimum width=4.55cm,minimum height=1.10cm,
    inner sep=2pt,anchor=center,align=center] (test) at (0,4.64)
    {$x^2-4x+3\leq 0$};
  \node[flowbox,minimum width=2.00cm] (update) at (0,3.22) {$x=x+1$};
  \node[flowbox,minimum width=2.20cm] (inc) at (0,1.90) {$n=n+1$};
  \node[io,minimum width=2.00cm] (output) at (3.20,3.22) {输出 $n$};
  \node[draw,rounded corners=7pt,minimum width=1.35cm,minimum height=.70cm]
    (finish) at (3.20,1.90) {结束};

  % Main chain and the single merge before the decision.
  \draw[flowarrow] (start.south) -- (input.north);
  \draw[flowarrow] (input.south) -- (init.north);
  \draw (init.south) -- (merge);
  \draw[flowarrow] (merge) -- (test.north);

  % 是 branch and its loop return.
  \draw[flowarrow] (test.south) -- (update.north);
  \draw[flowarrow] (update.south) -- (inc.north);
  \coordinate (loop-bottom) at (0,0.70);
  \coordinate (loop-left) at (-3.55,0.70);
  \coordinate (loop-entry) at (-3.55,5.43);
  \draw (inc.south) -- (loop-bottom) -- (loop-left) -- (loop-entry);
  \draw[flowarrow] (loop-entry) -- (merge);

  % 否 branch to output and finish.
  \draw[flowarrow] (test.east) -- (3.20,4.64) -- (3.20,3.62);
  \draw[flowarrow] (output.south) -- (finish.north);

  \node[inner sep=0pt,anchor=west] at (0.22,4.00) {是};
  \node[inner sep=0pt,anchor=west] at (2.45,5.00) {否};

  \path[use as bounding box] (-3.90,8.95) rectangle (3.65,0.35);
\end{tikzpicture}
